\documentclass[a4paper,12pt]{article}
\usepackage{color}
\usepackage{graphicx, gensymb}
\usepackage{amsfonts, cancel, amssymb}
\usepackage{amsmath, amsthm, array, esvect, esint}
\usepackage{typearea}
\usepackage[a4paper,left=2cm,right=2cm,top=2cm,bottom=3cm,bindingoffset=5mm]{geometry}
\newcommand\numberthis{\addtocounter{equation}{1}\tag{\theequation}}
\newcommand{\bra}{}
\newcommand{\kom}{}
\newcommand{\brak}{}
\newcommand{\braket}{}
\newcommand{\intx}{}
\newcommand{\intr}{}
\newcommand{\kla}{}
\newcommand{\klam}{}
\newcommand{\klammer}{}
\newcommand{\oper}{}
\newcommand{\tecxt}{}
\newcommand{\Bra}{}
\newcommand{\Ket}{}

\begin{document}

\title{07 Drehimpulsalgebra}
\author{Joachim Schwardt}
\date{\today}
\maketitle

\section{Drehimpulsoperator}
Mit $\nabla=e_r\partial_r+e_\vartheta\frac{1}{r}\partial_\vartheta+e_\varphi\frac{1}{r\sin\vartheta}\partial_\varphi$ folgt:
\begin{align*}
L&=r\times p=re_r\times\frac{\hbar}{i}\nabla=\frac{\hbar}{i}\kla\left( {e_\varphi\partial_\vartheta - e_\vartheta\frac{1}{\sin\vartheta}\partial_\varphi } \right)=\frac{\hbar}{i}\begin{pmatrix}
-\sin\varphi\partial_\vartheta - \cot\vartheta\cos\varphi\partial_\varphi \\ \cos\varphi\partial_\vartheta - \cot\vartheta\sin\varphi\partial_\varphi \\ \partial_\varphi
\end{pmatrix}\\
L^2&=-\hbar^2\kla\left( {e_\varphi\partial_\vartheta - e_\vartheta\frac{1}{\sin\vartheta}\partial_\varphi } \right)\cdot\kla\left( {e_\varphi\partial_\vartheta - e_\vartheta\frac{1}{\sin\vartheta}\partial_\varphi } \right)=-\hbar^2\kla\left( {\partial_\vartheta^2+\frac{\cos\vartheta}{\sin\vartheta}\partial_\vartheta+\frac{1}{\sin^2\vartheta}\partial_\varphi^2} \right) \\
L_\pm &=L_x\pm iL_y=\hbar\kla\left( {i\sin\varphi\partial_\varphi + i\cos\varphi\cot\vartheta\partial_\varphi\pm \cos\varphi\partial_\vartheta \mp \sin\varphi\cot\vartheta\partial_\varphi} \right)= \\
&=\hbar\kla\left( {\pm e^{\pm i\varphi}\partial_\vartheta +ie^{\pm i\varphi}\cot\vartheta\partial_\varphi} \right)=\hbar e^{\pm i\varphi}\kla\left( {i\cot\vartheta\partial_\varphi\pm\partial_\vartheta} \right) \\
\nabla^2&=\kla\left( {e_r\partial_r+e_\vartheta\frac{1}{r}\partial_\vartheta+e_\varphi\frac{1}{r\sin\vartheta}\partial_\varphi} \right)\cdot\kla\left( {e_r\partial_r+e_\vartheta\frac{1}{r}\partial_\vartheta+e_\varphi\frac{1}{r\sin\vartheta}\partial_\varphi} \right)=\\
&=\partial_r^2+\frac{1}{r^2}\partial_\vartheta^2+\frac{1}{r^2\sin^2\vartheta}\partial_\varphi^2+\frac{2}{r}\partial_r+\frac{\cos\vartheta}{r^2\sin\vartheta}\partial_\vartheta=\frac{1}{r^2}\partial_r\kla\left( {r^2\partial_r} \right)-\frac{1}{\hbar^2r^2}L^2
\end{align*}
\section{Kommutatoren}
Es sei $|m\rangle$ ein Eigenzustand zu $L_z$, d.h. $L_z|m\rangle = m\hbar |m\rangle$ und $\langle m|L_z=\langle m|m\hbar$. Sei $(j\in\klammer\left\{ {x,y} \right\})$:
\begin{align*}
\bra\langle {L_j}\rangle &=\frac{\epsilon_{kzj}}{i\hbar}\langle \left[ {L_k,L_z} \right]\rangle=\frac{\epsilon_{kzj}}{i\hbar}\braket\langle {m} | {L_kL_z-L_zL_k} | {m}\rangle=\frac{\epsilon_{kzj}}{i\hbar}\braket\langle {m} | {L_k m\hbar - m\hbar L_k} | {m}\rangle=0
\end{align*}
Es ist $L_\pm L_\mp=L_x^2+L_y^2\pm i[L_y,L_x]=L^2-L_z^2\pm\hbar L_z$.
\begin{align*}
\bra\langle {L_x^2-L_y^2}\rangle&=\frac{1}{i\hbar}\bra\langle {L_x[L_y,L_z]+[L_x,L_z]L_y}\rangle=\frac{1}{i\hbar}\bra\langle {[L_xL_y,L_z]}\rangle=\\
&=\frac{1}{i\hbar}\braket\langle {m} | {L_xL_ym\hbar-m\hbar L_xL_y} | {m}\rangle=0
\end{align*}
Aus $\bra\langle {L_x^2}\rangle=\bra\langle {L_y^2}\rangle$ folgt:
\begin{align*}
\bra\langle {L_x^2}\rangle&=\bra\langle {L^2}\rangle-\bra\langle {L_z^2}\rangle-\bra\langle {L_y^2}\rangle=\hbar^2j(j+1)-\hbar^2m^2-\bra\langle {L_x^2}\rangle\ \ \Rightarrow\ \ \bra\langle {L_x^2}\rangle=\frac{\hbar^2}{2}\kla\left( {j(j+1)-m^2} \right)
\end{align*} 
\section{Unschärferelation}
Es gilt:
\begin{align*}
\Delta L_j^2&=\bra\langle {(L_j-\bra\langle {L_j}\rangle)^2}\rangle=\bra\langle {L_j^2}\rangle=\frac{\hbar^2}{2}\kla\left( {j(j+1)-m^2} \right)\\
\Delta L_x\Delta L_y&=\frac{\hbar^2}{2}\kla\left( {j(j+1)-m^2} \right)\ge\frac{\hbar^2}{2}\ \ \ \ (\text{minimal für }m=\pm j)
\end{align*}
Sei $e=e_x\cos\alpha +e_y\cos\beta +e_z\cos\gamma$:
\begin{align*}
\bra\langle {e\cdot L}\rangle&=\bra\langle \cos\gamma L_z\rangle=m\hbar\cos\gamma \\
(\Delta e\cdot L)^2&=\bra\langle {(e\cdot L)^2}\rangle-\bra\langle {e\cdot L}\rangle^2=(\cos^2\alpha+\cos^2\beta)\bra\langle {L_x^2}\rangle+\cos^2\gamma\bra\langle {L_z^2}\rangle-m^2\hbar^2\cos^2\gamma= \\
&=\frac{\hbar^2}{2}(\cos^2\alpha+\cos^2\beta)\kla\left( {j(j+1)-m^2} \right)
\end{align*}
\section{Zweidimensionaler \textsc{Hilbert}-Raum}
Matrixdarstellung:
\begin{align*}
\sigma_x&=\begin{pmatrix}
0&1\\ 1&0
\end{pmatrix}=\sigma_x^\dagger &\sigma_y&=\begin{pmatrix}
0&-i\\ i&0
\end{pmatrix}=\sigma_y^\dagger &\sigma_z &=\begin{pmatrix}
1&0\\ 0&-1
\end{pmatrix}=\sigma_z^\dagger
\end{align*}
Es gilt für das charakteristische Polynom:
\begin{align*}
\chi_{\sigma_j}(\lambda)&=\lambda^2-\lambda\mathrm{tr\,}\sigma_j+\det\sigma_j=\lambda^2-1\ \ \Rightarrow\ \ \lambda=\pm 1
\end{align*}
Die Eigenwerte sind also für alle drei Matrizen gleich. Die Eigenvektoren berechnen sich wie folgt:
\begin{align*}
\sigma_xv_x&=\pm v_x\ \ \Rightarrow\ \ \left\{\begin{matrix}
0x + 1y = \pm x \\ 1x + 0y = \pm y
\end{matrix}\right.\ \ \Rightarrow\ \ v_x=\begin{pmatrix}
1 \\ \pm 1
\end{pmatrix} \\
\sigma_yv_y&=\pm v_y\ \ \Rightarrow\ \ \left\{\begin{matrix}
0x - iy = \pm x \\ ix + 0y = \pm y
\end{matrix}\right.\ \ \Rightarrow\ \ v_y=\begin{pmatrix}
1 \\ \pm i
\end{pmatrix} \\
\sigma_zv_z&=\pm v_z\ \ \Rightarrow\ \ \left\{\begin{matrix}
1x + 0y = \pm x \\ 0x - 1y = \pm y
\end{matrix}\right.\ \ \Rightarrow\ \ v_z=\frac{1}{2}\begin{pmatrix}
1\pm 1 \\  1\mp 1
\end{pmatrix}
\end{align*}
Mehr im Sinne der Dirac-Schreibweise und der Struktur der Quantenmechanik ist allerdings folgender Lösungsweg über $\sigma_j\Ket| {\lambda_j}\rangle =\lambda_j\Ket| {\lambda_j}\rangle $:
\begin{align*}
\braket\langle {n} | {\sigma_x} | {m}\rangle\brak\langle {m}| {\lambda_x}\rangle&=\lambda_x\brak\langle {n}| {\lambda_x}\rangle\\
\Rightarrow\ \ \brak\langle {-}| {\lambda_x}\rangle&=\lambda_x\brak\langle {+}| {\lambda_x}\rangle\tecxt\text{ \ und \ }\brak\langle {+}| {\lambda_x}\rangle=\lambda_x\brak\langle {-}| {\lambda_x}\rangle\\
\Rightarrow\ \ \lambda_x&=\pm 1\ \ \Rightarrow\ \ \Ket| {\lambda_x^\pm}\rangle = \frac{1}{\sqrt{2}}\kla\left( {\Ket| {+}\rangle \pm\Ket| {-}\rangle } \right)\\
\braket\langle {n} | {\sigma_y} | {m}\rangle\brak\langle {m}| {\lambda_y}\rangle&=\lambda_y\brak\langle {n}| {\lambda_y}\rangle\\
\Rightarrow\ \ -i\brak\langle {-}| {\lambda_y}\rangle&=\lambda_y\brak\langle {+}| {\lambda_y}\rangle\tecxt\text{ \ und \ }i\brak\langle {+}| {\lambda_y}\rangle=\lambda_y\brak\langle {-}| {\lambda_y}\rangle\\
\Rightarrow\ \ \lambda_y&=\pm 1\ \ \Rightarrow\ \ \Ket| {\lambda_y^\pm}\rangle = \frac{1}{\sqrt{2}}\kla\left( {\Ket| {+}\rangle \pm i\Ket| {-}\rangle } \right)\\
\braket\langle {n} | {\sigma_z} | {m}\rangle\brak\langle {m}| {\lambda_z}\rangle&=\lambda_z\brak\langle {n}| {\lambda_z}\rangle\\
\Rightarrow\ \ \brak\langle {+}| {\lambda_z}\rangle&=\lambda_z\brak\langle {+}| {\lambda_z}\rangle\tecxt\text{ \ und \ }\brak\langle {-}| {\lambda_z}\rangle=\lambda_z\brak\langle {-}| {\lambda_z}\rangle\\
\Rightarrow\ \ \lambda_z&=\pm 1\ \ \Rightarrow\ \ \Ket| {\lambda_z^\pm}\rangle = \Ket | {\pm} \rangle
\end{align*}
Man kann die Paul-Matrizen auch in anderen Basen darstellen. Berechne hierzu folgende Terme (Summen-Konvention beachten):
\begin{align*}
\sigma_j^q&=\braket\langle {\lambda_{jk}^p} | {\sigma_j^p} | {\lambda_{jl}^p}\rangle=\brak\langle {\lambda_{jk}^q}| {n}\rangle\braket\langle {n} | {\sigma_j^p} | {m}\rangle\brak\langle {m}| {\lambda_{jl}^q}\rangle\\
\sigma_x^x&=\frac{1}{\sqrt{2}}\begin{pmatrix}
1 & 1 \\ 1 & -1
\end{pmatrix}^\dagger \begin{pmatrix}
0 & 1 \\ 1 & 0
\end{pmatrix}\frac{1}{\sqrt{2}}\begin{pmatrix}
1 & 1 \\ 1 & -1
\end{pmatrix}= \begin{pmatrix}
1 & 0 \\ 0 & -1
\end{pmatrix}=\sigma_z^z\\
\sigma_y^x&=\frac{1}{\sqrt{2}}\begin{pmatrix}
1 & 1 \\ 1 & -1
\end{pmatrix}^\dagger \begin{pmatrix}
0 & -i \\ i & 0
\end{pmatrix}\frac{1}{\sqrt{2}}\begin{pmatrix}
1 & 1 \\ 1 & -1
\end{pmatrix}= \begin{pmatrix}
0 & i \\ -i & 0
\end{pmatrix}=-\sigma_y^z\\
\sigma_z^x&=\frac{1}{\sqrt{2}}\begin{pmatrix}
1 & 1 \\ 1 & -1
\end{pmatrix}^\dagger \begin{pmatrix}
1 & 0 \\ 0 & -1
\end{pmatrix}\frac{1}{\sqrt{2}}\begin{pmatrix}
1 & 1 \\ 1 & -1
\end{pmatrix}= \begin{pmatrix}
0 & 1 \\ 1 & 0
\end{pmatrix}=\sigma_x^z\\
\sigma_x^y&=\frac{1}{\sqrt{2}}\begin{pmatrix}
1 & 1 \\ i & -i
\end{pmatrix}^\dagger \begin{pmatrix}
0 & 1 \\ 1 & 0
\end{pmatrix}\frac{1}{\sqrt{2}}\begin{pmatrix}
1 & 1 \\ i & -i
\end{pmatrix}= \begin{pmatrix}
0 & -i \\ i & 0
\end{pmatrix}=\sigma_y^z\\
\sigma_y^y&=\frac{1}{\sqrt{2}}\begin{pmatrix}
1 & 1 \\ i & -i
\end{pmatrix}^\dagger \begin{pmatrix}
0 & -i \\ i & 0
\end{pmatrix}\frac{1}{\sqrt{2}}\begin{pmatrix}
1 & 1 \\ i & -i
\end{pmatrix}= \begin{pmatrix}
1 & 0 \\ 0 & -1
\end{pmatrix}=\sigma_z^z\\
\sigma_z^y&=\frac{1}{\sqrt{2}}\begin{pmatrix}
1 & 1 \\ i & -i
\end{pmatrix}^\dagger \begin{pmatrix}
1 & 0 \\ 0 & -1
\end{pmatrix}\frac{1}{\sqrt{2}}\begin{pmatrix}
1 & 1 \\ i & -i
\end{pmatrix}= \begin{pmatrix}
0 & 1 \\ 1 & 0
\end{pmatrix}=\sigma_x^z
\end{align*}
Es gilt $\sigma_j^2=I$. Betrachte außerdem $\sigma_i\sigma_j$:
\begin{align*}
\sigma_x\sigma_y&=\begin{pmatrix}
i & 0 \\ 0 & -i
\end{pmatrix}=i\sigma_z &\sigma_y\sigma_x&=\begin{pmatrix}
-i & 0 \\ 0 & i
\end{pmatrix}=-i\sigma_z \\
\sigma_y\sigma_z&=\begin{pmatrix}
0 & i \\ i & 0
\end{pmatrix}=i\sigma_x &\sigma_z\sigma_y&=\begin{pmatrix}
0 & -i \\ -i & 0
\end{pmatrix}=-i\sigma_x \\
\sigma_z\sigma_x&=\begin{pmatrix}
0 & 1 \\ -1 & 0
\end{pmatrix}=i\sigma_y &\sigma_x\sigma_z&=\begin{pmatrix}
0 & -1 \\ 1 & 0
\end{pmatrix}=-i\sigma_y
\end{align*}
Zusammengefasst gilt also:
\begin{align*}
\sigma_i\sigma_j&=\delta_{ij}I + i\epsilon_{ijk}\sigma_k
\end{align*}
Daraus kann man die Kommutator- und Antikommutatorrelation herleiten:
\begin{align*}
[\sigma_i,\sigma_j]&=\delta_{ij}I+i\epsilon_{ijk}\sigma_k-\delta_{ji}I-i\epsilon_{jik}\sigma_k=2i\epsilon_{ijk}\sigma_k \\
\klammer\left\{ {\sigma_i,\sigma_j} \right\}&=\delta_{ij}I+i\epsilon_{ijk}\sigma_k+\delta_{ji}I+i\epsilon_{jik}\sigma_k=2\delta_{ij}I
\end{align*}
\end{document}